\subsection {Untersuchungsgebiet und verwendete Geodaten}

Der in dieser Arbeit vorgestellte Workflow wurde für die urbane Vegetation der Stadt Leipzig entwickelt. Aufgrund der hohen urbanen Vegetationsdichte und der heterogenen Struktur, bestehend aus Parkanlagen, einzelnen Straßenbäumen sowie dem Leipziger Auwald, bietet die Stadt ein anspruchsvolles und repräsentatives Untersuchungsgebiet. Für die detaillierte Analyse wurden gezielt die Gohliser Straße und der Mariannenpark ausgewählt. Diese Gebiete spiegeln unterschiedliche Bestandsdichten urbaner Vegetation wider, wodurch die Übertragbarkeit des automatisierten Workflows auf das gesamte Stadtgebiet gewährleistet wird.

Die primäre Datengrundlage bilden ALS-LiDAR-Punktwolken, welche die erforderlichen 3D-Informationen liefern. Diese Daten wurden vom Landesamt für Geobasisinformation Sachsen (\ac{GeoSN}) bereitgestellt und stammen aus einer Befliegung vom 09. Januar 2023. Da die Erfassung im Winter stattfand, ist bei der Auswertung eine saisonal bedingt geringere Laubdichte sowie ein niedrigerer NDVI zu berücksichtigen. Die Punktwolken liegen im komprimierten LAZ-Format vor und weisen eine Punktdichte von mindestens vier Punkten pro $\text{m}^2$ auf. Die geometrische Genauigkeit ist mit einer Höhengenauigkeit von bis zu 0,15 m (Bezugssystem: Deutsches Haupthöhennetz 2016) und einer Lagegenauigkeit von 0,30 m spezifiziert. Die Daten sind bereits vom GeoSN vorklassifiziert in Bodenpunkte (Klasse 02) und Nichtbodenpunkte (Klasse 20). Zur Schließung von Messlücken wurden die Punktwolke zudem durch interpolierte Gewässerpunkte (Klasse 08) sowie interpolierte sonstige Punkte (Klasse 30), beispielsweise unterhalb von Gebäuden, ergänzt. Das GeoSN stellt diese Daten in Kacheln von 2 km x 2 km zum Download bereit \cite{zotero-item-43}.

Zusätzlich werden Digitale Orthofotos (\ac{DOP}) des GeoSN mit einer Bodenauflösung von 0,20 m genutzt, die vom 09. März 2025 stammen. Diese zeitliche Verschiebung zu den LiDAR-Daten impliziert potenzielle geometrische Abweichungen. Die DOPs liegen als Drei- und Vierkanalfarbbild vor. Das RGB-DOP (Rot, Grün, Blau) dient der visuellen Echtfarbendarstellung, während das RGBI-DOP (Rot, Grün, Blau, Nahes Infrarot/NIR) die Berechnung des NDVI ermöglicht. Dieser Index ist essenziell für die spektrale Filterung nach Vegetation sowie für die semantische Informationsanreicherung der Punktwolke \cite{zotero-item-43}.

Um die konsistente Verschneidung der Datensätze zu gewährleisten, werden alle Daten einheitlich im amtlichen Koordinatenreferenzsystem EPSG:25833 (ETRS89/UTM Zone 33N) verarbeitet. Dieses System kombiniert ein physikalisch definiertes geodätisches Bezugssystem für die Abbildung der Erdoberfläche mit einer konformen Zylinderprojektion für das ebene Lagekoordinatensystem \cite{zotero-item-43}.

Ergänzend zu den Daten des Landesamtes werden auch kommunale Datenbestände der Stadt Leipzig herangezogen. Als zentrale Referenz dient das städtische Baumkataster, das als punktbasierter Vektordatensatz im GeoPackage-Format (\ac{GPKG}) vorliegt. Für jedes enthaltene Baumobjekt sind semantische Attribute wie die Baumgattung, das Pflanzjahr sowie Schätzwerte zu Kronendurchmesser und Baumhöhe hinterlegt \cite{zotero-item-46}.
 